% Faithful editable transcription of selected portions of the published CC BY 4.0 PDF.
% Source: Flynn, Ramos, Young, Algebraic Combinatorics 9(4) (2026), 973--992, DOI 10.5802/alco.498.
% Physical PDF page 1 is the license cover; physical page k corresponds to printed page k+971 for k>=2.
% BEGIN context partitions -- printed p.976, physical PDF p.5
Throughout this work we will write $\mathfrak{S}_n$ to denote the symmetric group on $n$-letters.
\begin{sourceblock}{Definition 2.1}
Let $n\geq 0$ be fixed. A \textbf{partition} of $n$ is a tuple $\lambda=(\lambda_1,\ldots,\lambda_r)$ of positive integers such that $\sum_i\lambda_i=n$, and $\lambda_i\geq\lambda_{i+1}$ for each $i$. By convention, when $n=0$, we say that there is precisely one partition of $n$, namely the empty partition.

We write $\mathbb{Y}_n$ to denote the set of all partitions of the integer $n$.

Relatedly, a \textbf{partition of $[n]$} is a collection of non-empty, disjoint subsets of the set $[n]=\{1,\ldots,n\}$ whose union is all of $n$. If $\lambda$ is some partition of $n$, we say that a partition of $[n]$, $\alpha$, has shape $\lambda$ if the sizes of the subsets present in $\alpha$ can be organized to form $\lambda$. We write $\mathbb{P}_n$ to denote the set of all partitions of $[n]$, and $\mathbb{P}_n(\lambda)$ to denote the subset of partitions with shape $\lambda$.
\end{sourceblock}
It is a well known fact that the irreducible complex representations of $\mathfrak{S}_n$ are in bijection with partitions of $n$. We will generally write $S^\lambda$ to denote the irreducible representation associated to the partition $\lambda$.

\textbf{It is our convention that the irreducible corresponding to the partition $(n)$ is the trivial representation.}

\begin{sourceblockroman}{Remark 2.2}
It is common to visualize partitions of $n$ as tableau of boxes, where the $i$-th row of boxes has $\lambda_i$ boxes. This visualization is referred to as the \textbf{Young diagram} associated to $\lambda$. We will frequently jump between this visualization and the ``tuple-form'' of a partition for certain definitions. For instance, in Definition 2.10, we define one partition to be the transpose of another.
\end{sourceblockroman}
% END context partitions

% BEGIN context tabloids -- printed p.977, physical PDF p.6
\begin{sourceblockroman}{Remark 2.3}
While it is not absolutely critical to know precisely how the irreducible representation partition correspondence works, there is one particular presentation of these modules that we will continuous refer to, and so we take a moment to define it now. Fix a partition $\lambda$ of $n$. Consider the set $\mathcal{T}$ of all bijections $T$ from the Young diagram of $\lambda$ to $[n]$. That is to say, $T$ assigns in a bijective fashion numbers $\{1,2,\ldots,n\}$ to each box of the Young diagram. If $R(T)\subseteq\mathfrak{S}_n$ is the subgroup of permutations which preserve the elements in the rows of $T$, then we write $\{T\}$ to denote the orbit of $T$ under the action of $R(T)$. The object $\{T\}$ is often referred to as a \textbf{tabloid} with shape $\lambda$.

The set of all tabloids of shape $\lambda$ form an $\mathfrak{S}_n$-set in the obvious way, and we may write $M^\lambda$ to denote the $\mathfrak{S}_n$-representation which linearizes this set. The irreducible representation $S^\lambda$ can be found as a subrepresentation of $M^\lambda$ as follows. For each bijection $T$ from the diagram of $\lambda$ to $[n]$, define
\[
 v_T=\sum_{\sigma\in C(T)}\operatorname{sgn}(\sigma)\{\sigma(T)\},
\]
where $C(T)$ is the subgroup of permutations that preserve the elements in the columns of $T$. The irreducible representation $S^\lambda$ is the span of the vectors $v_T$.
\end{sourceblockroman}
% END context tabloids

% BEGIN context rationality -- printed p.977, physical PDF p.6
\begin{sourceblock}{Theorem 2.5}
Let $n\geq 0$ be an integer, and let $\lambda$ be a partition of $n$. Then $\mathfrak{S}_n$ acts on the vector space $S^\lambda$ by matrices with rational entries. In particular, the dual of any complex symmetric group representation is isomorphic to the original representation.
\end{sourceblock}
% END context rationality

% BEGIN context padding -- printed p.977, physical PDF p.6
\begin{sourceblock}{Definition 2.6}
Let $\lambda$ be a fixed partition of some integer, and let $n\geq |\lambda|+\lambda_1$. Then we define the \textbf{padded partition}
\[
 \lambda[n]:=(n-|\lambda|,\lambda)
\]
We will typically write $S_n^\lambda$ for the irreducible representation corresponding to $\lambda[n]$. Moreover, we will implicitly assume that $S_n^\lambda=0$ for any $n<|\lambda|+\lambda_1$.
\end{sourceblock}
% END context padding

% BEGIN context FI -- printed pp.977--978, physical PDF pp.6--7
\begin{sourceblock}{Definition 2.7}
Let $\mathrm{FI}$ denote the category whose objects are the sets $[n]=\{1,\ldots,n\}$, and whose morphisms are injections of sets. An \textbf{FI-module over a ring $k$} is a functor,
\[
 V:\mathrm{FI}\to k-\mathrm{mod},
\]
from $\mathrm{FI}$ to the category of finitely generated $k$-modules. We will usually write $V_n:=V([n])$ and $f_*:=V(f)$.

We say that an FI-module $V$ is \textbf{finitely generated} if there is an integer $N$ such that for all $n\geq N$, the group $V_n$ is generated as an $\mathfrak{S}_n$-module by the image of $V_N$ under the map induced by the standard inclusion $[N]\hookrightarrow[n]$ In this case we also say that $V$ is \textbf{generated in degrees $\leq N$}.
\end{sourceblock}
% END context FI

% BEGIN context FI stability -- printed p.978, physical PDF p.7
\begin{sourceblock}{Theorem 2.8 ([4])}
Let $V$ be a finitely generated FI-module over a field $k$ of characteristic $0$ that is generated in degrees $\leq N$. Then,
\begin{enumerate}
\item There is a finite set $\Lambda\subseteq\bigcup_n\mathbb{Y}_n$, as well as non-negative integers $\{m_\lambda\}_{\lambda\in\Lambda}$, such that for all $n\gg0$,
\[
 V_n\cong\bigoplus_{\lambda\in\Lambda}m_\lambda S_n^\lambda;
\]
\item The set $\Lambda$ from the above part consists exclusively of partitions with at most $N$ parts;
\item There exists a polynomial $p_V\in\mathbb{Q}[x_1,x_2,\ldots]$ of degree $\leq N$, such that for all $n\gg0$, the character $\chi_n$ of $V_n$ satisfies,
\[
 \chi_n(\sigma)=p_V(X_1(\sigma),X_2(\sigma),\ldots),
\]
where $X_i(\sigma)$ is the number of $i$-cycles in the decomposition of $\sigma$. In particular, there is a polynomial $p\in\mathbb{Q}[n]$ such that for all $n\gg0$, $\dim_{\mathbb{Q}}V_n=p_V(n)$.
\end{enumerate}
\end{sourceblock}
% END context FI stability

% BEGIN context intrinsic -- printed pp.978--979, physical PDF pp.7--8
\textbf{For the remainder of this section we fix a partition $\lambda$ of some non-negative integer $n$.}
\begin{sourceblock}{Definition 2.10}
Let $\lambda^\dagger$ denote the \textbf{conjugate partition} of $\lambda$, obtained from $\lambda$ by transposing its Young diagram. If $\alpha\in\mathbb{P}_n(\lambda^\dagger)$ and $\alpha$ has parts $\alpha_1,\alpha_2,\ldots,\alpha_r$ such that $[n]=\cup_i\alpha_i$, we define $\mathfrak{S}_\alpha$ to be the subgroup $\prod_i\mathfrak{S}_{\alpha_i}$ of $\mathfrak{S}_n$. We will also write $T_\alpha$ to denote a set of transpositions that generate the group $\mathfrak{S}_\alpha$.

With notation as above, we may now define the \textbf{intrinsic hyperplane arrangement} of $S^\lambda$ to be that comprised of the hyperplanes
\[
 H_\alpha:=\sum_{\tau\in T_\alpha}(S^\lambda)^\tau,
\]
where $(S^\lambda)^\tau$ invariant subspace of $S^\lambda$ with respect to the transposition $\tau$. We write $\mathcal{A}_\lambda$ for this hyperplane arrangement.
\end{sourceblock}
It is entirely non-obvious that the above definition is well defined. However, it is shown in [23] that the spaces $H_\alpha$ are hyperplanes, and that they do not depend on the choice of generating transpositions $T_\alpha$.

One should observe that for any permutation $\sigma\in\mathfrak{S}_n$, the action of $\sigma$ on $S^\lambda$ descends to an action of $\sigma$ on $\mathcal{A}_\lambda$. Indeed, it isn't hard to see that
\[
 \sigma\cdot H_\alpha=H_{\sigma\cdot\alpha}
\]
for any $\alpha\in\mathbb{P}_n(\lambda^\dagger)$. Indeed, this just follows from the fact that for any transposition $\tau$, $v\in S^\lambda$ is fixed by $\tau$ if and only if $\sigma\cdot v$ is fixed by $\sigma\tau\sigma^{-1}$. Later we will see that this action on the hyperplane arrangement can be lifted to an action by the category $\mathrm{FI}$.
% END context intrinsic

% BEGIN context complement -- printed pp.979--980, physical PDF pp.8--9
\begin{sourceblock}{Definition 2.12}
Let $(V,\mathcal{A})$ denote a complex hyperplane arrangement\footnote{In this work all hyperplane arrangements are assumed to be central, and finite.} and for each $H\in\mathcal{A}$ choose an orthogonal vector $a(H)$. The \textbf{complement of $\mathcal{A}$} is then defined to be the variety
\[
 \mathcal{M}(\mathcal{A})=\{v\in V\mid\langle v,a(H)\rangle\neq0\text{ for all }H\in\mathcal{A}\}
\]
The variety $\mathcal{M}(\mathcal{A})$ can be embedded into the torus $(\mathbb{C}^\times)^{\mathcal{A}}$ by setting the $H$-th coordinate to be $\langle a(H),\bullet\rangle$. We define $\mathcal{M}^{-1}(\mathcal{A})$ to be the image of $\mathcal{M}(\mathcal{A})$ under the involution of the torus which inverts every coordinate. Finally, the \textbf{reciprocal plane} $X_{\mathcal{A}}$ is then defined to be the closure of $\mathcal{M}^{-1}(\mathcal{A})$ in the affine space $\mathbb{C}^{\mathcal{A}}$.

The \textbf{Kazhdan-Lusztig polynomial} of the hyperplane arrangement $\mathcal{A}$, $p_{\mathcal{A}}(t)$, is defined to be the Poincar\'e polynomial associated to the intersection cohomology. That is to say, the coefficient in front of $t^i$ in $p_{\mathcal{A}}(t)$ is precisely given by $\dim_{\mathbb{C}}IH^{2i}(X_{\mathcal{A}};\mathbb{C})$.

If the arrangement $\mathcal{A}$ is equipped with an action by some finite group $\Gamma$, then this action descends to an action on intersection homology. In this case we may define the \textbf{equivariant Kazhdan-Lusztig polynomial} $p_{\mathcal{A}}^\Gamma(t)$ to be the formal polynomial whose coefficient of $t^i$ is the virtual representation associated with $IH^{2i}(X_{\mathcal{A}};\mathbb{C})$.
\end{sourceblock}
% END context complement

% BEGIN context arrangements -- printed p.980, physical PDF p.9
\begin{sourceblock}{Definition 3.1}
Let $(V,\mathcal{A})$ and $(W,\mathcal{B})$ be linear hyperplane arrangements over $\mathbb{C}$. A \textbf{morphism} between them $\varphi:(V,\mathcal{A})\to(W,\mathcal{B})$ is a surjection of vector spaces $\varphi^*:W\to V$ such that for any $H\in\mathcal{A}$, one has that $(\varphi^*)^{-1}(H)\in\mathcal{B}$. In this way, a morphism of arrangements induces both:
\begin{itemize}
\item A regular map $\mathcal{M}(\mathcal{B})\to\mathcal{M}(\mathcal{A})$, where $\mathcal{M}$ indicates the complement of the hyperplane arrangement and;
\item A map between the semi-lattices induced by the intersections of the hyperplanes defining $\mathcal{A}$ and $\mathcal{B}$, respectively. We write $\mathcal{L}(\mathcal{A})$ to denote the intersection semi-lattice of $\mathcal{A}$. We usually will refer to this as the \textbf{lattice of flats} of $\mathcal{A}$
\end{itemize}
An FI-arrangement is a functor $(V_\bullet,\mathcal{A}_\bullet)$ from $\mathrm{FI}$ to the category of arrangements. We say that an FI-arrangement $(V_\bullet,\mathcal{A}_\bullet)$ is \textbf{finitely generated} if there is an integer $N$ such that for all $n\geq N$ and every subspace $S\in\mathcal{L}(\mathcal{A}_n)$, $S$ can be written as an intersection of preimages of subspaces coming from $\mathcal{L}(\mathcal{A}_N)$. That is to say, for any $S\in\mathcal{L}(\mathcal{A}_n)$ one can find $S_1,\ldots,S_r\in\mathcal{L}(\mathcal{A}_N)$ as well as injections $f_i:[N]\to[n]$, such that
\[
 S=\cap_i(V(f_i)^*)^{-1}(S_i),
\]
Finally, an FI-arrangement is \textbf{normal} if for every morphism $f:[m]\hookrightarrow[n]$ in $\mathrm{FI}$ and every subspace $S\in\mathcal{L}(\mathcal{A}_n)$ that contains the kernel of the map induced by $f$ on vector spaces, $S$ is in the image of the map induced by $f$ on the semi-lattice of intersections
\end{sourceblock}
\begin{sourceblockroman}{Remark 3.2}
Note that for a given map between arrangements $\varphi:(V,\mathcal{A})\to(W,\mathcal{B})$, the corresponding map between vector spaces and between hyperplane complements goes in the opposite direction, whereas the corresponding map between semi-lattices goes in the same direction. In particular, an FI-arrangement does not in general immediately induce an FI-module by composing with the forgetful map to vector spaces.
\end{sourceblockroman}
% END context arrangements

% BEGIN context Gadish -- printed p.981, physical PDF p.10
\begin{sourceblock}{Theorem 3.4 ([9, Theorem A])}
If $(V_\bullet,\mathcal{A}_\bullet)$ is a finitely generated and normal FI-arrangement, then for any $i\geq0$ the rational cohomology of the complement space,
\[
 H^i(\mathcal{M}(\mathcal{A}_\bullet);\mathbb{Q})
\]
forms a finitely generated FI-module over $\mathbb{Q}$.
\end{sourceblock}
% END context Gadish

% BEGIN proof construction -- printed p.981, physical PDF p.10
Our way forward has therefore been made clear. We must show,
\begin{itemize}
\item for any partition $\lambda$, the collection of hyperplane arrangements $(S_n^\lambda,\mathcal{A}_{\lambda[n]})$ can be glued together in some natural way to form an FI-arrangement and,
\item that this FI-arrangement is both finitely generated and normal.
\end{itemize}
We take on these two problems in turn.

To begin, fix a partition $\lambda$ and recall from [4] that there exists a (finitely generated) FI-module over $\mathbb{Q}$ whose evaluation at $[n]$ is the module $S_n^\lambda$ so long as $n\geq\lambda_1+|\lambda|$ and the zero module otherwise. Denote this FI module by $S_\bullet^\lambda$. By dualizing, one obtains an $\mathrm{FI}^{\mathrm{op}}$-module whose evaluation at $[n]$ is the module
\[
 (S_n^\lambda)^\vee\cong S_n^\lambda,
\]
by Theorem 2.5. In fact, for any injection $f:[m]\hookrightarrow[n]$, the maps induced by $f$ on the original FI-module, $f_*$, and its dual, $f^*$, are partial inverses of one another. This grants us that the map induced by $f:[m]\to[n]$,
\[
 f^*:S_n^\lambda\to S_m^\lambda
\]
is surjective, and that for any $\alpha\in\mathbb{P}_m(\lambda[m]^\dagger)$
\[
 (f^*)^{-1}(H_\alpha)=H_{f(\alpha)},\text{ where }f(\alpha)\in\mathbb{P}_n(\lambda[n]^\dagger)
\]
In particular, we have an FI-arrangement whose evaluation at each $n$ is precisely the intrinsic hyperplane arrangement,
\[
 (S_n^\lambda,\mathcal{A}_{\lambda[n]}),
\]
as desired. To prove our representation stability theorem it therefore remains to show that this FI-arrangement is both finitely generated and normal.
% END proof construction

% BEGIN proof normality -- printed pp.981--982, physical PDF pp.10--11
\begin{sourceblock}{Proposition 3.5}
The FI-arrangement $(S_\bullet^\lambda,\mathcal{A}_{\lambda[\bullet]})$ is both finitely generated and normal.
\end{sourceblock}
\begin{proof}
To begin, we will continue to use the convention that a set partition $\alpha\in\mathbb{P}_n(\lambda[n]^\dagger)$ is identified with the pieces of the partition of size strictly greater than $1$. In this way $\mathbb{P}_n(\lambda[n]^\dagger)$ becomes an $\mathfrak{S}_n$ set by the obvious action. In fact, one may think of $\mathbb{P}_n(\lambda[n]^\dagger)$ as being a (finitely generated) FI-set in the sense of [19]. In particular, for any $n\gg0$ and any finite collection of elements in $\mathbb{P}_n(\lambda[n]^\dagger)$, $\{\alpha_i\}_i$, one may find injections $f_i:[n-1]\hookrightarrow[n]$ and partitions $\{\beta_i\}_i$ such that for each $i$, $\alpha_i=f_i(\beta_i)$. Finite generation now follows from this as well as the fact that the action of $\mathrm{FI}$ on $\mathcal{A}_{\lambda[\bullet]}$ agrees with its action on $\mathbb{P}_\bullet(\lambda[\bullet]^\dagger)$.

It therefore remains to show that these arrangements are normal. Let $f:[m]\hookrightarrow[n]$ be an injection of sets. We first hope to understand which subspaces in $\mathcal{L}(\mathcal{A}_{\lambda[n]})$ will contain the kernel of the map induced by $f$ on the underlying vector spaces. More specifically, we only need to understand when a given hyperplane $H_\alpha$ contains the kernel of this map, for $\alpha\in\mathbb{P}_n(\lambda[n]^\dagger)$. Equivalently, we need to understand when the cokernel of the forward map induced by $f$,
\[
 f_*:S_m^\lambda\to S_n^\lambda,
\]
is contained in $H_\alpha$. From the tabloid description of Remark 2.3, it can be seen that
\[
 f_*(v_T)=v_{f(T)},
\]
where $f(T)$ is obtained from $T$ by first post-composing with $f$, and then filling in the extra padding in the first row with the elements outside of the image of $f$. In particular, the cokernel of the map induced in the forward direction by $f$ will be spanned by those vectors $v_T$ for which at least one element of $[n]$ outside the image of $f$ appears outside of the extra padding in the first row. It follows that $H_\alpha$ can only possibly contain this cokernel if the parts of $\alpha$ exclusively contain elements from the image of $f$. Indeed, otherwise, the normal vector to $H_\alpha$, as described in [23], is among the vectors $v_T$ that span the cokernel.

Now given some intersection of hyperplanes with this condition on $\alpha$, We find that it will be in the image of the map induced by $f$ on the intersection semi-lattice by looking at the intersection of hyperplanes whose associated set partition contains the relevant preimages (under $f$) of those image values that appeared in the original hyperplanes' associated partition.
\end{proof}
% END proof normality

% BEGIN statement main -- printed p.983, physical PDF p.12
\begin{sourceblock}{Theorem 3.7}
Fix a partition $\lambda$ and $i\geq0$, and write $V=H^i(\mathcal{M}(\mathcal{A}_{\lambda[\bullet]});\mathbb{Q})$. Then,
\begin{enumerate}
\item There is a finite set $\Lambda\subseteq\bigcup_n\mathbb{Y}_n$, as well as non-negative integers $\{m_\mu\}_{\mu\in\Lambda}$, such that for all $n\gg0$,
\[
 V_n\cong\bigoplus_{\mu\in\Lambda}m_\mu S_n^\mu;
\]
\item The set $\Lambda$ from the above part consists exclusively of partitions with at most $i*(|\lambda|+\lambda_1)$ parts;
\item There exists a polynomial $p_V\in\mathbb{Q}[x_1,x_2,\ldots]$ of degree $\leq i\cdot(|\lambda|+\lambda_1)$, such that for all $n\gg0$, the character $\chi_n$ of $V_n$ satisfies,
\[
 \chi_n(\sigma)=p_V(X_1(\sigma),X_2(\sigma),\ldots),
\]
where $X_i(\sigma)$ is the number of $i$-cycles in the decomposition of $\sigma$. In particular, there is a polynomial $p\in\mathbb{Q}[n]$ such that for all $n\gg0$, $\dim_{\mathbb{Q}}V_n=p_V(n)$.
\end{enumerate}
\end{sourceblock}
% END statement main

% BEGIN proof main -- printed p.983, physical PDF p.12
\begin{proof}
In view of Proposition 3.5, as well as Theorems 3.4 and 2.8, the only part of this theorem that remains unjustified are the bounds of $i\cdot(|\lambda|+\lambda_1)$ that appear throughout. We will accomplish this using a few key facts from the study of hyperplane arrangements. Firstly, the famous theorem of Orlik and Solomon [15] tells us that for any fixed $n\geq|\lambda|+\lambda_1$,
\begin{itemize}
\item the first cohomology group $H^1(\mathcal{M}(\mathcal{A}_{\lambda[n]});\mathbb{Q})$ is a vector space with basis in bijection with the hyperplanes of $\mathcal{A}_{\lambda[n]}$, and
\item for any $i>1$, the cohomology group $H^i(\mathcal{M}(\mathcal{A}_{\lambda[n]});\mathbb{Q})$ is a quotient of the tensor power $(H^1(\mathcal{M}(\mathcal{A}_{\lambda[n]});\mathbb{Q}))^{\otimes i}$
\end{itemize}
The first fact can be applied to show that, as an FI-module, $H^1(\mathcal{M}(\mathcal{A}_{\lambda[\bullet]});\mathbb{Q})$ is generated in degree $|\lambda|+\lambda_1$. The bounds of the original Theorem 2.8 now imply our bounds for this module. For higher $i$, one simply applies the fact that generating degree is additive in tensor powers of FI-modules [4], and the second fact found above. Note that while this is not enough to prove that the higher cohomologies are generated in degrees at most $i\cdot(|\lambda|+\lambda_1)$, it certainly is strong enough to conclude the upper bounds we are looking for, as upper bounds on the irreducible decomposition and dimension growth will descend to quotients.
\end{proof}
% END proof main

% BEGIN context method scope -- printed p.983, physical PDF p.12
\begin{sourceblockroman}{Remark 3.8}
One might observe that the above proof can be enhanced so-as to entirely avoid the theory of FI-arrangements. Indeed, it would essentially amount to showing that the quotient maps from the tensor powers of the first cohomology to the $i$-th cohomology respected all of the FI-structure. Such a proof, however, would essentially only give us statements about these cohomology groups at the level of algebra. In contrast, our method can be thought of as proving finite generation at the level of the topological spaces themselves. This has utility in-so-far-as applications to \'etale cohomology, as detailed in [9].
\end{sourceblockroman}
% END context method scope
